\documentclass[11pt]{article}
\usepackage{syntonly}
\usepackage{verbatim}
\usepackage{graphicx}%input graphs.
\usepackage{mdwlist}
\usepackage{amsmath} %tag and eqref
\usepackage{amsthm}
\usepackage{textcomp} %makes the degree symbol also available
\usepackage{minted}
%\textdegree \textcelsius
\hyphenation{FORTRAN}
\pagestyle{plain}
\title{\LaTeX{}  learning}
\author{suchan}
\begin{document}
\maketitle
\verbatiminput{a.html}
\begin{figure}
	\centering
	\scalebox{1.0}{
	\includegraphics{test}
}
\caption{sdfsdf}

\end{figure}
\inputminted[linenos]{c}{avg_realloc.c}
``This'' 'is' `my' \emph{first} document prepared in" \LaTeX.\\
\textbackslash documentclass\{article\} tells \LaTeX that what we want to produce is an article. If you want to write a book, this must be changed to \textbackslash documentclass\{book\}. The whole document we want to typeset should be included between \textbackslash begin\{document\} and \textbackslash end\{document\}.
\\* Then \textbackslash emph\{first\}, becomes\emph{first} in the output. It is a common practice to \emph{emphasize} words in print using italic letters. I typed it on \today.

\section{Input files}
\subsection{space}
It does not matter whether you enter one or several     spaces after a word.

An empty line starts a new paragraph.
\subsection{Special Characters}
The following characters are reserved characters that either have a special meaning under \LaTeX or are not available in all font.If you type them directly in \LaTeX, they will normally not print.

As you can see.... \# \$ \% \^{} \& \_ \{ \} \~{} \textbackslash...\^{} and \~{} should be \textbackslash \^{}\{\} and \textbackslash \~{}\{\}.
\subsection{\LaTeX{} Commands}
\LaTeX ignore whitespace after commands. If you want to get a space after a command, you have to put either an empty parameter{} and a blank or special spacing command after the command name. The empty parameter\{\} stops \LaTeX{} from eating up all whitespace after the command name.

Some commands require a perameter, which has to be given between curly braces\{ \} after the command name. Some commands take optional parameters, which are inserted after the command name in square brcakets\[ \].

\textbackslash \emph{Command} [\emph{Optional parameter} ] \{\emph{parameter}\}

Some examples:

You can \textsl{lean} on me! You can \textbackslash testsl \{lean\} on me.

Please, start a new line right here!\textbackslash newline \newline
Thankyou!

\subsection{Comments}

\%

The \% character can also be used to split long input lines where no whitespace or line breaks\%are allowed.

For longer comments, you can use verbatim package. Add line \textbackslash usepackage\{verbatim\} to the preamble of your document.

This is another \begin{comment} abcdefg \end{comment}

This is another \textbackslash begin\{comment\} abcdefg \textbackslash end \{comment\}

\subsection{Input File Structure}
Every inputfile start with the command \textbackslash documentclass\{...\}

To load a package use the command: \textbackslash usepackage\{...\}

You start with the body of the text with the command:\textbackslash begin\{document\}

At the end of the document you add \textbackslash end\{document\}

\subsection{A Typical Command Line Session}
The area between \textbackslash documentclass and \textbackslash\{document\} is called a preamble.

Compiling your \LaTeX{} input file:
\begin{comment}
	\begin{enumerate}
	\begin{itemize}
	\begin{description} \hfill
	\item Nested
\end{comment}
\begin{description}
	\item[first] \hfill \\ Create your \LaTeX input file. This file must be plain ASCII text.
	\item[second] \hfill  
		\begin{enumerate}
		\item
		Open shell.
	\item
		latex foo.tex
	\item
		xdvi foo.dvi \&
\end{enumerate}
\end{description}
Convert dvi file to PostScript for printing or viewing with GhostScript.
		
dvips -Pcmz foo.dvi -o foo.ps

dvipdf foo.dvi

\subsection{The Layout of the Document}
\subsubsection{Document Classes}
Example:An input file for a \LaTeX{} document could start with\linebreak[3]\break \sloppy the line.

\textbackslash documentclass[11pt, twoside,a4paper]\{article\}

\subsubsection{Packages}
If you want to include graphics, coloured text or source code from a file into your document, you need to enhance the capability of \LaTeX
. Such enhancement called packages.

Add commandline: \textbackslash usepackage[options]\{package\}
\begin{itemize}
	\item
		article for articles in scientific journals, short reports, presentations, invitations, program documentations.
	\item
		proc a cladd for proceeding based on article class.
	\item
		minimal is as small as it can get. It only set page size and base font. It is mainly used for debugging purpose.
	\item
		report for longer reports containing several chapters, small books, PHD thesis,...
	\item
		book for real books
	\item
		slide for slides. 
	\item
\end{itemize}


\section{typesetting}
\subsection{Line breaking and page breaking}
For a line breaking, using \textbackslash\textbackslash or \textbackslash newline
The difference between ``the empty line between two sentences" and \textbackslash\textbackslash is intended.
Starting a new page: using \textbackslash newpage.
It seems that \textbackslash linebreak[n] and \textbackslash pagebreak[n] can only be used above 4..because you leave \LaTeX the option of ignoring your command if the result would look very bad.
\subsection{Hyphenation}
This means that if you place a hyphenation command into the preamble of your document it will influence th e English language hyphenation. If you place the command after the \textbackslash begin\{document\} and you are using some package for national language support like babel, then the hyphenation hints will be active in the language activated through babel.

\paragraph{If you use \textbackslash\{FORTRAN Hy-phen-a-tion\} can prevents FORTRAN, no special character or symbols are allowed in the argument.}

command \textbackslash - inserts a discretionary hyphen into a word. I think this isaaaaaaaaaaaaaaaaaaaaaaaaaaaaaaaaaaa:su\-persdf\-sdfsdf\-gggggg.

Several worda can be kept together on one line with the command by using \textbackslash mbox{}and \textbackslash fbox{} \mbox{12345678912312312} \fbox{12123123123123\\sdhjfsdh}

\subsection{Ready-made String}
\today

\LaTeX

\TeX

\LaTeXe

\subsection{Special Characters and Symbols}
\subsubsection{Quotation Marks}
You should not use " in \LaTeX{} , use two grave accent for opening and the vertical quote for closing quotation.
\subsubsection{Dashes and Hyphens}
X-rated

pages 12--37

yes---or no?

$0$, $1$ or $-1$
\subsubsection{Tilde($\sim$)}
\~{} and $\sim$.

\subsubsection{Slash}
read\slash write allows hyphenation. Normal `/' character may be still used for ratios or units. 5MB/s.

\subsubsection{Degree Symbol($^{\circ}\mathrm{C}$)}

30 \textcelsius{}
45 \textdegree{}F.
\texteuro%use package{textcomp}

\subsubsection{Ellipsis}
\ldots(low dots)
Not like this ... but like this:\\
New York. Tokyo, Budapest, \ldots

\subsubsection{Ligatures}
Disable ligatures: shelf\mbox{}ful

se\~no h\^otel \"o
I like BASIC\@.~ What about you

sdfsdf\label{sec:this} sdfsd~\ref{sec:this}
sdhfjsdhf\footnote{ssdfsd} \underline{text}
\emph{a\emph{c}}c
\verb@\ldots@
\subsubsection{tabular}
\begin{tabular}{|r|l|}
	\hline
	aaa&bbb\\
	ccc&ddd\\
	eee&ffff\\ \cline{2-2}
	ggg&hhh\\ \hline
	\hline
	rrr&uuu\\
	\hline
\end{tabular}
\begin{tabular}{|p{4.7cm}|}
	\hline
	aaaaaaaaaaaaaaa\\
	\hline
\end{tabular}

\begin{tabular}[c]{|p{10cm}|c|r|l|}
	\hline
	111111&222222&333333&444444\\
	\hline
	111111&222222&333333&444444\\
	\hline
\end{tabular}

\begin{tabular}{@{dddd} c @{fffff}}
	\hline
	sdfsdfsdf\\
	\hline
\end{tabular}

\begin{tabular}{|c|c|c|}
	\hline
	\multicolumn{3}{|c|}{aaa}\\
	\hline
	wwww&yyyyyyyyyyyy&yyuiu\\
	\hline
\end{tabular}
\begin{tabular}{c r @{.} l}
	pi &\multicolumn{2}{c}{value}\\
	\hline
	$\pi$ & 3&1415\\
	$\pi^{\pi}$ &  45&6\\
\end{tabular}

\begin{tabular}{|p{4.7cm}|}
\hline
Add a squared and b squared to get c squared.
\begin{equation}
	a^2+b^2=c^3
\end{equation}
Einstein says:
\begin{equation}
	E=mc^2\label{clever}
\end{equation}
He didn't say
\begin{equation}
	1+1=3 \tag{dumb}
\end{equation}
This is a reference to \eqref{clever}.\\
\hline
\end{tabular}

In text style:
$1\frac{1}{2} hours$
$1\frac{1}{2}$~hours

$1\dfrac{1}{2} hours$ \qquad
$1\tfrac{1}{2} hours$

$\sqrt{\frac{x^2}{y+1}}$
$\surd{\frac{x^2}{y+1}}$
$x^{\frac{2}{k+1}}$
$\frac{\partial^2f}{\partial{x^2}}$

$\binom{n}{k}=\binom{n-1}{k}+\binom{n-1}
{k-1}$ 
$f_n(x)\stackrel{*}{\approx}1 \, sdfjkslsk?$
\begin{equation}
\sum_{i=1}^n
\prod_\epsilon
\end{equation}
sdfsfdsdf
sdfsdfsfd

sdfsdf

\verb.sdfsdf.


\begin{proof}
	thisidjf
	\[e^{x+y}=e^x\cdots e^y \qedhere\]
\end{proof}
\begin{proof}
	thisidjf
	$$e^{x+y}=e^x\cdots e^y \qedhere$$
\end{proof}

sdfsdf
\end{document}
